\documentclass[../main.tex]{subfiles}
\begin{document}
\section{Introduction}
\label{sec:introduction}
%---background---
Empirical regression coefficients can depend strongly on relatively small subsets of the estimation sample. Classical case-deletion diagnostics provide a direct way to study this dependence one observation at a time \parencite{cook2000detection,belsley2005regression}. Their observation-level structure can become less informative when influence is distributed across several observations, particularly under masking and related forms of joint influence \parencite{rousseeuw1990unmasking,hadi1993procedures}. Recent data-dropping methods have renewed attention to this broader sensitivity problem by asking how empirical conclusions change when selected groups of observations are removed \parencite{broderick2020automatic,kuschnig2021hidden}.

Outlyingness and target-coefficient influence are related but distinct properties. An observation can be unusual because of its outcome, covariates, or relationship to the fitted model while exerting limited influence on a prespecified regression coefficient. Conversely, substantial coefficient movement can arise from a group of observations that is not conventionally classified as an outlier set. This thesis refers descriptively to \emph{concentrated target-coefficient sensitivity} when deletion of a relatively small subset produces substantial movement in the coefficient of interest. The relevant question is therefore not only which observations are unusual, but how coefficient sensitivity becomes distributed and concentrated across the estimation sample.

Most Influential Sets (MIS) provide a coefficient-targeted framework for studying this structure. For a prespecified deletion budget \(k\), fixed-\(k\) MIS searches jointly for the size-\(k\) set that maximises movement in a target regression coefficient \parencite{kuschnig2021hidden,konrad2026finding}. The resulting set characterises joint influence under a stated coefficient target, fitted specification, and deletion budget. This perspective differs from robust estimation, which modifies the estimation rule to obtain coefficient stability under contamination, and from specification testing, which evaluates particular implications of model departures. These distinctions make it possible to study localisation, estimation, and model misspecification as related but separate aspects of coefficient sensitivity.

The central research question is: \emph{How do data contamination and model misspecification generate concentrated sensitivity in a prespecified regression coefficient, and what does fixed-\(k\) MIS reveal about the structure and consequences of that sensitivity?} Three simulation studies address complementary parts of this question. The first examines when response outliers, good leverage, and bad leverage align conventional outlyingness with joint target-coefficient influence. The second examines how influential-set localisation relates to post-deletion estimation and to robust estimators designed for coefficient stability. The third examines whether model misspecification can generate concentrated coefficient sensitivity without a planted contamination set, and whether the selected observations correspond to substantively defined affected subgroups.

The simulations show that the structure of coefficient sensitivity is strongly mechanism dependent. Bad leverage creates close alignment between planted contamination and target-slope influence, while response outliers and good leverage produce weaker correspondence between unusualness and coefficient movement. Localisation and coefficient accuracy also need not coincide: identifying the set that maximises coefficient movement addresses a different statistical objective from obtaining a stable estimator under contamination. The misspecification results extend this distinction beyond planted outliers. Several model departures generate strong concentrated sensitivity, while the fitted-span structure of the linear-endogeneity design produces an invariance boundary in which severity changes the coefficient representation without changing the standardised deletion-sensitivity structure.

The empirical analysis examines the magnitude and structure of the same sensitivity across ten published studies built around randomized experimental variation. One prespecified estimand from each study is evaluated under a common deletion protocol with full-refit validation. The audits show substantial cross-study heterogeneity in the deletion fraction required to produce large coefficient movements. They also reveal different influential-set structures, including persistent influential cores, non-nested membership paths across deletion budgets, and geographic concentration. These results characterise finite-sample dependence of the audited coefficients within their validated specifications.

The contribution of the thesis is therefore to use fixed-\(k\) influential-set analysis to organise several dimensions of target-coefficient sensitivity within a common framework. The analysis distinguishes unusualness from joint coefficient influence, localisation from robust estimation, and concentrated sensitivity from model correction. It also connects these distinctions to empirical patterns in realised estimation samples. The deletion budget \(k\) remains a sensitivity parameter, and the selected set is specific to the target coefficient and maintained specification.

The remainder of the thesis proceeds as follows. Section 2 reviews classical influence diagnostics, joint influence, data-dropping methods, robust regression, and model misspecification. Section 3 defines the fixed-\(k\) MIS methodology, including the set-DFBETA, linear-fractional optimisation, statistical calibration, and the role of the deletion budget. Section 4 presents the simulation evidence on outlyingness, influence localisation and estimation, and misspecification-induced coefficient sensitivity. Section 5 reports cross-study empirical evidence on the magnitude and structure of target-coefficient sensitivity. Section 6 concludes and discusses the scope and future development of influential-set analysis.

\end{document}
